\section{Envelopes with a fixed power gain}\label{sec:envelopes}

We construct envelopes whose gradient integral grows geometrically under
iteration.  The part of each parent cube available for the next iteration
will have more gradient integral than the entire parent had before the
change.  All scales and gain factors are fixed before any finite
approximation order is prescribed.

Throughout this Section, tensor norms in Hessian inequalities are operator
norms for $G$.  Coordinate derivatives and integrals refer to fixed smooth
coordinates.  The notation $p,Q,J_h$ is that of
\eqref{eq:hessian-test}.

\subsection{A seed and the preparation condition}

\begin{lemma}\label{lem:seed}
There are a number $a_*>0$, a smooth real function $h_0$ on $B$, a
$C^1$ neighborhood $\mathcal G$ of $g_0$ among smooth metrics, and a
coordinate chart $\Omega\subset\{t>1/2\}$ with open sets
$U_0\Subset V\Subset\Omega$ such that the following assertions hold
uniformly for $G\in\mathcal G$:
\begin{enumerate}[label=\textup{(\roman*)}]
 \item $h_0=-a_*t$ on $\{t\le1/2\}$;
 \item $\max_e J_{h_0}(e)\ge c_0>0$ on $\{t\ge1/4\}$;
 \item on $U_0$, the Hessian of $h_0$ is positive definite with a
 uniform positive lower bound, and
 $0<s_{\min}\le|\nabla_Gh_0|_G\le s_{\max}$.
\end{enumerate}
The maximum in \textup{(ii)} is over all admissible tests, including all
unit vectors when $\dd h_0=0$.
\end{lemma}

\begin{proof}
Begin with $\widetilde h=-a_*t$.  Direct differentiation on the round
sphere gives
\begin{equation}\label{eq:seed-identities}
 |\nabla_{g_0}\widetilde h|_{g_0}^2
 =4a_*^2t(1-t),\qquad
 \nabla_{g_0}^2\widetilde h
 =2a_*(tg_0-\dd y\cdot\dd y).
\end{equation}
At every point there is a nonzero tangent vector in $\ker\dd y$,
because $\dd y:T_xB\to\R^3$ has a four-dimensional domain.  A round
unit vector $e$ in this kernel satisfies $\dd\widetilde h(e)=0$ and
$(\nabla_{g_0}^2\widetilde h)(e,e)=2a_*t$.  Since the Hessian in
\eqref{eq:seed-identities} has norm at most $2a_*$ and the squared
gradient has norm at most $a_*^2$, on $t\ge1/4$ we obtain
\[
 J_{\widetilde h}(e)\ge a_*/2-2a_*^3\ge a_*/4
\]
after choosing $a_*^2\le1/8$.

At a point $x_*$ with $t=1$, the Hessian has eigenvalues
$2a_*,2a_*,0,0$, and the gradient is zero.  On a small neighborhood
of $x_*$ contained in $\{t>1/2\}$, two Hessian eigenvalues therefore
have a uniform positive lower bound and the gradient is arbitrarily
small.  In round normal coordinates $\xi$ at $x_*$, add
$\varepsilon|\xi|^2/2$ times a fixed smooth cutoff equal to one near
$\xi=0$.  Choose $\varepsilon>0$ sufficiently small that this is a
$C^2$-small perturbation on its support.  The resulting function $h_0$
has positive definite Hessian at $x_*$.

Here is why the cutoff region retains the required test.  Shrinking its
support and then the perturbation if necessary, the perturbed Hessian
has two eigenvalues at least $c a_*$, its norm is at most $C a_*$,
and $|\nabla_{g_0}h_0|$ is sufficiently small that
$C a_*|\nabla_{g_0}h_0|^2<c a_*/2$.  The positive two-plane
intersects $\ker\dd h_0$ in dimension at least one.  A unit vector
in that intersection gives $J_{h_0}\ge c a_*/2$.  This argument also
applies at a critical point.  Outside the perturbation support the
previous $\ker\dd y$ test applies.

These conclusions persist in a sufficiently small common $C^1$
neighborhood of $g_0$.  Indeed, the difference of the two Hessians of
a fixed function is the connection difference contracted with its
differential, and all metric and gradient comparisons use only
$C^1$ bounds.  Outside the perturbation support, the condition
$\dd y(e)=0$ still implies $\dd h_0(e)=0$ independently of the
metric.  On the support, the positive two-plane and small-gradient
estimates just established are uniform under a small perturbation.

Finally, positive definiteness at $x_*$ holds on a neighborhood of that
point.  There is a nearby point with nonzero differential: otherwise
$h_0$ would be constant on that neighborhood.  Choose a relatively
compact neighborhood $U_0$ of such a point where the Hessian remains
positive definite and the gradient stays separated from zero.  Choose
$U_0\Subset V\Subset\Omega\subset\{t>1/2\}$ in a fixed
coordinate chart, and shrink $\mathcal G$ once more.  All the stated
constants can then be chosen uniformly.  Notice that $U_0$ is chosen
away from the critical point $x_*$.
\end{proof}

We use the coordinates of $\Omega$ for all corrugations.  They may be
taken on a neighborhood of its closure, so metric and inverse metric
have common ellipticity and $C^1$ bounds there.  Fix positive constants
$C_{\mathrm p}$ and $c_{\mathrm p}$.  A smooth function $h$ is
\emph{prepared at scale $S$} on a cube $P$ of side $\nu S$,
$0<S\le1$, if
\begin{equation}\label{eq:preparation}
 \begin{gathered}
 |\nabla h|>0,\qquad
 |\nabla^2h|\le C_{\mathrm p}|\nabla h|/S,\\
 \text{at each point there is a two-plane }
 E\subset\ker\dd h\text{ such that }
 J_h(e)\ge c_{\mathrm p}|\nabla h|/S
 \quad(e\in E,\ |e|=1).
 \end{gathered}
\end{equation}
The plane $E$ need not vary continuously.  Cubes are parallel to the
fixed coordinate axes, and the inequalities are required on their
closures.

The two Hessian conditions serve different parts of the next corrugation.
The Hessian norm bound controls the inherited Hessian and the variation
of the parent gradient.  When the gradient changes by a small relative
amount, we can rotate the favorable two-plane slightly into the
hyperplane perpendicular to the new gradient.  One further linear
tangency constraint then leaves a direction with a positive test.
On the region selected for repetition, the modification supplies two
favorable directions again.  This is why the preparation condition
asks for a two-plane even though the envelope ultimately needs only one
positive test at each point.

By Lemma~\ref{lem:seed}, for all sufficiently large
$C_{\mathrm p}$ and sufficiently small $c_{\mathrm p}>0$, every
sufficiently small cube in $U_0$ is prepared for $h_0$ at scale $1$,
uniformly over $G\in\mathcal G$.

\subsection{A repeatable local corrugation}

\begin{lemma}\label{lem:corrugation}
Under fixed ellipticity and $C^1$ metric bounds in the coordinate
patch, there are constants $C_{\mathrm p}^*>0$,
$c_{\mathrm p}^*>0$ and $\zeta>0$ with the following property.
For any $C_{\mathrm p}\ge C_{\mathrm p}^*$ and
$0<c_{\mathrm p}\le c_{\mathrm p}^*$, one may choose a sufficiently
small reciprocal integer $\nu>0$, once and for all, such that every
prepared cube $P$ at scale $S\le1$ admits a smooth modification
$v=h+w$, with $w$ supported in the interior of $P$, satisfying:
\begin{enumerate}[label=\textup{(\roman*)}]
 \item for fixed constants $0<\alpha\le1\le\beta$ and $c>0$,
 \[
  \alpha|\nabla h|\le|\nabla v|\le\beta|\nabla h|,
  \qquad
  \max_eJ_v(e)\ge c|\nabla h|/S
  \quad\hbox{on }P;
 \]
 \item writing $\ell=\nu^2S$, some cubes in the partition of $P$
 into cubes of side $\nu\ell$ are prepared for $v$ at scale $\ell$
 with the same constants in \eqref{eq:preparation}, and their union
 $\mathcal R_P$ satisfies
 \begin{equation}\label{eq:local-mass-gain}
  \int_{\mathcal R_P}|\nabla v|\,\dd x
  \ge(1+\zeta)\int_P|\nabla h|\,\dd x;
 \end{equation}
 \item if $s=|\nabla h(x_P)|$ at the center $x_P$ of $P$, then for
 every integer $j\ge0$,
 \begin{equation}\label{eq:corrugation-derivatives}
  \|D^jw\|_{L^\infty(P)}\le C_j s\ell^{1-j}.
 \end{equation}
\end{enumerate}
All constants are uniform over prepared functions and cubes and over
the stated class of metrics.  They may depend on the fixed pattern,
preparation constants, and final choice of $\nu$.  No derivative of
the input $h$ beyond its prepared second-derivative bound is used.
The pattern and $\zeta$ are chosen before the preparation constants
and the final $\nu$; increasing $C_{\mathrm p}$ or decreasing
$c_{\mathrm p}$ only requires taking $\nu$ smaller.
\end{lemma}

\begin{proof}
The pattern will depend on three variables transverse to the parent
gradient frozen at its center.  Radial wells in these variables have
two positive tangential Hessian directions; a small amplitude controls
their potentially negative radial contribution to $J_v$.  We first
choose the wells and smaller concentric balls that will supply the
repeatable gain.  We then place this pattern in the parent and verify
the Hessian test, including where the parent cutoff varies.

\paragraph{The periodic wells.}
Choose a smooth function $F$ on $[0,1]$, positive on $(0,1)$, equal
to $v$ near $v=0$, and flat at $v=1$.  Inside a unit cube in $\R^3$
take finitely many balls with pairwise disjoint closures and radii
$R_i>0$.  On a ball, at radius $r$ from its center, set
\[
 \psi(r)=-\int_r^{R_i}F(q/R_i)\,\dd q,\qquad
 f(r)=F(r/R_i),\qquad u_0(r)=f(r)/r.
\]
Extend $\psi$ by zero outside the balls and periodically to $\R^3$.
It is smooth: near a center it is a constant plus $r^2/(2R_i)$,
and it is flat at the outer boundary.  The extended value of $u_0$
at the center is $1/R_i$.  Write $\sigma$ for the radial unit vector.
Then
\begin{equation}\label{eq:well-identities}
 \nabla\psi=f\sigma,\qquad
 D^2\psi=f'\sigma\otimes\sigma
              +u_0(I-\sigma\otimes\sigma),\qquad
 f^2|f'|\le C u_0,\quad f^2\le C u_0.
\end{equation}
The first constant in these inequalities is independent of the radii,
because, for $v=r/R_i$,
\[
 \frac{f^2|f'|}{u_0}=vF(v)|F'(v)|.
\]
The second follows from $f^2/u_0=R_i vF(v)$ and the bounded radii.
Set $f,f',u_0$ equal to zero outside the balls.  These coefficients
are continuous, and
\begin{equation}\label{eq:well-vanishing}
 u_0=0\quad\Longrightarrow\quad f=f'=\psi=0.
\end{equation}
At a center $f'=u_0$, so any choice of $\sigma$ gives the same
Hessian in \eqref{eq:well-identities}; outside the balls choose
$\sigma$ arbitrarily as well.

Choose a small amplitude $a>0$.  Its smallness will use only $F$ and
the radius-independent first bound in \eqref{eq:well-identities}.
In particular it can be chosen before the packing.  The mean gradient
factor on any one well is
\[
 \mathcal W(a)=3\int_0^1v^2\sqrt{1+a^2F(v)^2}\,\dd v>1.
\]
Put $\eta_0=(\mathcal W(a)-1)/8$.  Choose the finite packing with
volume fraction $\varpi$ such that
\[
 \varpi\mathcal W(a)>1+6\eta_0.
\]
Such a packing exists.  To see this without requiring an infinite
packing in the construction, fill a fixed positive fraction of any
remaining open set by finitely many closed grid cubes lying inside
it, and put a ball with strictly smaller radius than the inradius
in each cube.  This captures a fixed positive fraction of the
remaining volume.  Finitely many repetitions make the uncovered
volume as small as desired, with all ball closures still disjoint.

Choose $\eta>0$ so small that, on the core balls
$r<(1-\eta)R_i$, the weighted integral per period satisfies
\[
 C_{\mathrm{core}}
 :=\sum_i\int_{r<(1-\eta)R_i}\sqrt{1+a^2f^2}\,\dd z'
 >\varpi\mathcal W(a)-\eta_0>1+5\eta_0.
\]
Choose an inner cube fraction $q\in(0,1)$ such that
\begin{equation}\label{eq:repeat-slack}
 qC_{\mathrm{core}}>C_{\mathrm{core}}-\eta_0>1+4\eta_0.
\end{equation}
Here $q$ is a four-dimensional volume fraction.  Fix a smooth cutoff
on the unit parent cube which equals one on the concentric inner
cube of fraction $q$, is supported in the parent interior, and is
flat wherever it is zero.  Products of one-dimensional plateau
cutoffs give such a function.  The wells, cores, and cutoff are now
fixed.  All their derivative bounds are finite, although no bound
uniform in derivative order is asserted.

\paragraph{Normalization on a prepared parent.}
Take affine coordinates $z=(z_1,z')$ centered at $x_P$, making the
metric Euclidean at the center and
$\nabla h(x_P)=s\partial_{z_1}$.  The linear maps and their inverses
have uniform bounds; the image of $P$ is a parallelepiped of uniformly
bounded distortion.  The coordinate identity relating $D(\dd h)$
to the covariant Hessian and the bounded connection gives
\[
 |D(\dd h)|\le C|\dd h|/S
\]
on $P$, since $S\le1$.  Integration along segments of length
$O(\nu S)$ yields relative upper and lower bounds for the gradient
and then
\begin{equation}\label{eq:parent-normalization}
 \dd h/s=\dd z_1+O(\nu),\qquad
 (G_{ij})_z=I+O(\nu),\qquad
 |\nabla h|/s=1+O(\nu).
\end{equation}
The constants here may depend on the preparation constants.  Only
the bounds already specified in the lemma enter them.

Rescale the fixed parent cutoff to $P$ and call it $\chi$.  Set
\begin{equation}\label{eq:corrugation-formula}
 v=h+s\ell a\chi\psi(z'/\ell),\qquad \ell=\nu^2S.
\end{equation}
In the following estimates the radial coefficients are evaluated at
$z'/\ell$.  A cutoff derivative at its parent scale is denoted
$D_{\mathrm{sc}}\chi$; thus an ordinary derivative of $\chi$ is
$(\nu S)^{-1}D_{\mathrm{sc}}\chi$.  The bounded affine changes only
alter the fixed bounds for these scaled derivatives.  Since
$\ell/(\nu S)=\nu$, Equation~\eqref{eq:corrugation-formula} gives
\begin{equation}\label{eq:new-gradient}
 \dd v/s=\dd z_1+a\chi f\sigma\cdot\dd z'+O(\nu).
\end{equation}
This proves the gradient comparison in \textup{(i)} for fixed
$\alpha,\beta$, when $\nu$ is sufficiently small.

\paragraph{The new Hessian contribution.}
Write
\[
 D=\frac{\nabla v}{\sqrt{1+|\nabla v|^2}},\qquad
 \mathcal F=\{e:\dd v(e)=0,\ \sigma\cdot e'=0\},
\]
where $e'$ denotes the static transverse components in the affine
coordinates.  The two defining linear conditions are independent
by \eqref{eq:new-gradient}, so $\mathcal F$ is two-dimensional.
For a metric unit vector $e\in\mathcal F$,
\begin{equation}\label{eq:test-components}
 |e'|^2\ge\tfrac12,\qquad
 |\sigma\cdot D'|\le C_0a\chi f+C\nu.
\end{equation}
Indeed $e_1=O(\nu)$ follows from the two linear conditions.  The
second estimate follows from \eqref{eq:new-gradient} and
$G^{ij}=\delta_{ij}+O(\nu)$, using a denominator at least a fixed
multiple of $s$.  We may take $C_0$ absolute by first taking $\nu$
small for the fixed data.  These arguments are uniform even when
$s$ is arbitrarily small or large.

Let $\Pi_\sigma$ be Euclidean projection onto $\sigma^\perp$ in
the transverse space.  The leading contribution of $v-h$ to its
Hessian test on $D,e$, divided by $s/\ell$, is
\[
 a\chi\left[
 f'(\sigma\cdot D')^2+
 u_0\bigl(|\Pi_\sigma D'|^2+|e'|^2\bigr)\right].
\]
By \eqref{eq:test-components}, this is at least
\[
 a\chi\left[
 \tfrac12u_0-|f'|\bigl(2C_0^2a^2\chi^2 f^2+C\nu^2\bigr)
 \right].
\]
At the initial choice of amplitude, impose
$2C_0^2a^2\sup_v vF(v)|F'(v)|\le1/4$.
The main term is then at least
$a\chi u_0/4-Ca\nu^2\chi|f'|$.

The mixed cutoff terms on $e,e$ vanish because
$\sigma\cdot e'=0$.  On $D,D$ their size is bounded by
\[
 C\nu f|D_{\mathrm{sc}}\chi|(\chi f+\nu).
\]
The pure cutoff terms have size at most
$C\nu^2|\psi||D_{\mathrm{sc}}^2\chi|$.
Finally, replacing the coordinate Hessian of $v-h$ by its covariant
Hessian gives an error of size
\[
 C\ell(\chi f+\nu)
\]
in the same units, because the connection is uniformly bounded.

We verify uniformity of these errors also where the favorable
coefficient vanishes.  Set $X=\chi u_0$.  On a fixed compact
parameter space, if continuous functions $X,b\ge0$ satisfy
$b=0$ on $\{X=0\}$, then
\begin{equation}\label{eq:vanishing-coefficient}
 \sup\frac{\nu^2b}{X+\nu^2}\longrightarrow0
 \qquad(\nu\longrightarrow0).
\end{equation}
For $X\ge\tau$ the quotient is at most
$\nu^2\sup b/\tau$; for $X<\tau$ it is at most
$\sup_{X<\tau}b$, which tends to zero as $\tau\downarrow0$.
The compact parameter space here consists of cutoff position,
period position, and bounded affine frames.  Taking its closure is
harmless.  Transformed cutoff derivatives are continuous on this
space and have uniform bounds.

Apply \eqref{eq:vanishing-coefficient} to
$b=\chi|f'|$, $b=f|D_{\mathrm{sc}}\chi|$, and
$b=|\psi||D_{\mathrm{sc}}^2\chi|$.  Their vanishing follows from
\eqref{eq:well-vanishing} and the flatness of $\chi$ on its zero
set.  The other mixed term satisfies
$\nu\chi f^2|D_{\mathrm{sc}}\chi|\le C\nu X$.
Moreover, $\chi f\le C\sqrt X$ and $\ell\le\nu^2$ give
\[
 \frac{\ell(\chi f+\nu)}{X+\nu^2}
 \le C\frac{\nu^2\sqrt X}{X+\nu^2}+\nu
 \le C'\nu.
\]
We have proved, uniformly in all prepared parents and allowed tests,
\begin{equation}\label{eq:corrugation-hessian-error}
 \frac{\ell}{s}\left[
  \nabla^2(v-h)(D,D)+\nabla^2(v-h)(e,e)\right]
 \ge c_a X-\varepsilon(\nu)(X+\nu^2),
 \qquad \varepsilon(\nu)\longrightarrow0,
\end{equation}
for a fixed $c_a>0$.  No positive lower bound on a nonzero value of
$\chi$ or $u_0$ has been used.

\paragraph{Retaining one direction everywhere.}
The old Hessian has norm $O(\nu^2)$ after division by $s/\ell$.
Thus the preceding estimate compares a favorable new contribution of
size $X$ with inherited terms of size $\nu^2$.  When $X$ dominates
$\nu^2$, the new favorable directions suffice.  When $X=O(\nu^2)$,
the gradients are close enough to carry an old favorable direction to
the new tangent space.  To obtain one uniform lower bound, consider
any sequence $\nu\to0$ of admissible configurations.
If $X/\nu^2\to\infty$ along a subsequence, the positive term in
\eqref{eq:corrugation-hessian-error} dominates the old Hessian in
every unit direction of $\mathcal F$.

If instead $X/\nu^2$ stays bounded, then
$\chi f\le C\sqrt X=O(\nu)$, so
$|\nabla v-\nabla h|\le\varepsilon_1|\nabla h|$ with
$\varepsilon_1=O(\nu)$.  This also makes their normalized gradients
close uniformly for unbounded slopes.  In fact, for
$T(p)=p/\sqrt{1+|p|^2}$ and $|q-p|\le\varepsilon_1|p|$,
$\varepsilon_1\le1/2$, the segment from $p$ to $q$ stays outside
the ball of radius $|p|/2$.  Therefore
\begin{equation}\label{eq:normalized-gradient-stability}
 |T(q)-T(p)|
 \le\frac{\varepsilon_1|p|}{\sqrt{1+|p|^2/4}}
 \le2\varepsilon_1.
\end{equation}
There is a metric orthogonal rotation $R=I+O(\varepsilon_1)$ taking
$\ker\dd h$ to $\ker\dd v$.  Rotate the old good two-plane $E$
in \eqref{eq:preparation} by $R$.  Inside the three-dimensional
space $\ker\dd v$, the planes $RE$ and $\mathcal F$ intersect in
dimension at least
\[
 2+2-3=1.
\]
For a unit vector $e$ in their intersection, compare the old test on
$p/Q,R^{-1}e$ with the old Hessian evaluated on $D,e$.
Equation~\eqref{eq:normalized-gradient-stability} shows that the
change costs $o(\nu^2)$ in the normalized units.  Thus the old
contribution is at least $\tfrac12c_{\mathrm p}\nu^2$ for small
$\nu$.  The new contribution in
\eqref{eq:corrugation-hessian-error} is bounded below by
$-o(\nu^2)$ in this bounded-ratio regime.

These alternatives prove the uniform margin
$c_{\mathrm p}\nu^2/8$ in the normalized units for some test when
$\nu$ is sufficiently small.  Indeed, a contrary
sequence has a subsequence on which $X/\nu^2$ is bounded or tends
to infinity, and both alternatives just excluded failure of that
margin.  Multiplication by $s/\ell$ and
\eqref{eq:parent-normalization} prove the Hessian assertion in
\textup{(i)}.  This is where base dimension four is used: in
dimension $d$ the corresponding intersection lower bound is $d-3$.

\paragraph{The region prepared for repetition.}
Inside the inner parent cube where $\chi=1$, keep the points for
which $z'/\ell$ lies in a core ball.  On this region $u_0$ has a
fixed positive lower bound.  Equations
\eqref{eq:corrugation-hessian-error} and
\eqref{eq:parent-normalization} imply, for every unit
$e\in\mathcal F$,
\[
 J_v(e)\ge c' s/\ell,
 \qquad
 |\nabla^2v|\le(C'+O(C_{\mathrm p}\nu^2))s/\ell.
\]
The leading constants depend only on the fixed pattern and the
uniform metric bounds.  Gradient comparison now permits fixed
thresholds $C_{\mathrm p}^*$ and $c_{\mathrm p}^*$ such that, for
$C_{\mathrm p}\ge C_{\mathrm p}^*$ and
$0<c_{\mathrm p}\le c_{\mathrm p}^*$, this region satisfies
\eqref{eq:preparation} at scale $\ell$ once $\nu$ is small.
In particular the old term $O(C_{\mathrm p}\nu^2)$ is absorbed
after $C_{\mathrm p}$ has been fixed.

Select all closed child cubes contained in this open repeat region.
Since $\nu$ is a reciprocal integer, the parent has an exact
partition into cubes of side $\nu\ell$.  These children satisfy the
preparation condition.  It remains to show that their total gradient
integral exceeds that of the entire old parent.  To evaluate their mass,
first tile the affine $z$ coordinates by period cells of side
$\ell$, using a dummy period in $z_1$.  The child grid uses the original
coordinates, whereas these period cells use the affine coordinates;
the two grids need not align.  The cells meeting the inner
parent boundary have relative volume
$O(\ell/(\nu S))=O(\nu)$, uniformly over the bounded affine
distortions.  The points lost by retaining only whole child cubes
lie in a layer of width $O(\nu\ell)$ along the core and inner
parent boundaries.  In each rescaled period the core boundaries
have finite area, so this gives another relative $O(\nu)$ loss.
The constants may involve the smallest radius in the finite fixed
packing.

On the repeat region,
\[
 |\nabla v|/s=\sqrt{1+a^2f^2}+O(\nu).
\]
The constant affine Jacobian cancels on dividing by the parent
volume.  Consequently
\[
 \frac1{s|P|}\int_{\mathcal R_P}|\nabla v|\,\dd x
 \ge qC_{\mathrm{core}}-C\nu,
 \qquad
 \frac1{s|P|}\int_P|\nabla h|\,\dd x\le1+C\nu.
\]
Choose the final $\nu$ small enough that both errors are at most
$\eta_0$.  By \eqref{eq:repeat-slack}, the ratio of new selected
mass to old entire parent mass is then greater than
\[
 \frac{1+3\eta_0}{1+\eta_0}
 =1+\frac{2\eta_0}{1+\eta_0}.
\]
Fix $\zeta=\eta_0/(1+\eta_0)>0$.  This proves
\eqref{eq:local-mass-gain}; taking $\nu$ smaller still preserves
this value of $\zeta$.

\paragraph{Derivatives and the order of choices.}
Differentiating \eqref{eq:corrugation-formula} in the fixed
coordinates gives \eqref{eq:corrugation-derivatives}.  In a term
with $q$ derivatives on $\chi$, the scale factor is
\[
 s\ell(\nu S)^{-q}\ell^{-(j-q)}
 =s\ell^{1-j}\nu^q\le s\ell^{1-j}.
\]
All derivatives of the pattern and all powers of the bounded affine
matrix enter the fixed constant $C_j$.  The center slope and the
affine matrix are constants in the expression being differentiated.

To summarize the choices within this proof, first fix $F$ and its
small amplitude $a$, then the finite packing, cores, and parent
cutoff with the slack in \eqref{eq:repeat-slack}.  Choose the
preparation constants next, large and small respectively, and
finally choose a reciprocal $\nu$ satisfying the local positivity,
repeat preparation, and mass-error requirements.  This proves the
stated uniformity and completes the lemma.
\end{proof}

\subsection{Iteration and a common exponent for every finite order}

\begin{proposition}\label{prop:envelope}
After fixing the data of Lemma~\ref{lem:seed} and possibly shrinking
$\mathcal G$, there are common constants $\gamma>0$ and $c>0$
with the following property.  For every smooth $G\in\mathcal G$
and all sufficiently large $k$, there is a smooth real envelope
$h=h_k$ on $B$ such that $h=-a_*t+c_k$ on $\{t\le1/2\}$ and
\begin{equation}\label{eq:envelope-bounds}
 \begin{gathered}
  \|h\|_{C^j}\le C_j k^{\delta(j+1)}
  \quad\text{for every fixed integer }j\ge0,\\
  \max_{\substack{|e|_G=1\\\dd h(e)=0}}J_h(e)
       \ge c k^{-\delta}\quad\text{on }\{t\ge1/4\},
  \qquad
  \int_V\sqrt{1+|\nabla_Gh|_G^2}\,\dd x\ge c k^\gamma.
 \end{gathered}
\end{equation}
Here $\delta=10^{-6}$, and the chart $\Omega$ and region
$V\Subset\Omega$ are fixed.  All corrugations are supported in
one fixed initial cube compactly contained in $V$.
The constants $c_k$ may be any bounded prescribed real numbers.
The exponent $\gamma$ and all construction parameters are
independent of derivative orders, subsequent accuracy requirements,
and round gluing collars.  Constants $C_j$ may depend on the order,
the fixed smooth metric and coordinate systems, and the bound for
$|c_k|$; no uniform bound as $j\to\infty$ is asserted.
\end{proposition}

\begin{proof}
Choose $C_{\mathrm p}$ large enough and $c_{\mathrm p}>0$ small
enough for both Lemma~\ref{lem:corrugation} and the seed preparation
following Lemma~\ref{lem:seed}.  Choose the final $\nu$ sufficiently
small for the former lemma and so that a fixed cube $P_0$ of side
$\nu$ has closure in $U_0$.  All these choices are uniform for the
metrics in $\mathcal G$.  Put $r=\nu^2$.

Start with the prepared family $\mathcal P_0=\{P_0\}$ for $h_0$.
At generation $i\ge1$, apply Lemma~\ref{lem:corrugation} on every
cube of $\mathcal P_{i-1}$, using parent scale $S_i=r^{i-1}$,
and call the resulting function $h_i$.  The selected child cubes
form $\mathcal P_i$; their sides are $\nu r^i$ and they are
prepared at scale $r^i$.  Every generation consists of finitely
many cubes with disjoint interiors.  Each perturbation is supported
strictly inside its parent, so each $h_i$ is smooth, including across
all cube boundaries.  Nothing outside $P_0$ changes.

Let
\[
 I_i=\sum_{P\in\mathcal P_i}\int_P|\nabla h_i|\,\dd x.
\]
Equation~\eqref{eq:local-mass-gain} gives
\begin{equation}\label{eq:iterated-mass}
 I_i\ge(1+\zeta)I_{i-1},\qquad
 \int_V\sqrt{1+|\nabla h_i|^2}\,\dd x
 \ge I_i\ge I_0(1+\zeta)^i.
\end{equation}
The initial mass $I_0$ has a positive uniform lower bound by
Lemma~\ref{lem:seed} and the fixed volume of $P_0$.

We give the derivative accounting for all orders with one growth
constant.  Let $\beta\ge1$ and $0<\alpha\le1$ be the gradient
comparison factors of Lemma~\ref{lem:corrugation}.  The largest
gradient after $i$ generations is at most a fixed multiple of
$\beta^i$.  On $P_0$ every point has gradient at least
$s_{\min}\alpha^i$, since a point is changed at most once per generation.
At generation $i$, the center slope of each parent is therefore
at most $C\beta^{i-1}$.  By
\eqref{eq:corrugation-derivatives}, disjointness within that
generation gives
\begin{equation}\label{eq:generation-derivatives}
 \|D^j(h_i-h_{i-1})\|_{L^\infty}
 \le C_j\beta^{i-1}r^{i(1-j)}.
\end{equation}
There is no factor counting cubes in this supremum norm estimate.
There is also no product of higher pattern derivative constants
over generations: the increments are added, and their individual
formulas involve frozen slopes and affine matrices.  Passing to
the fixed finite atlas only adds fixed chain-rule constants and
lower-order derivatives, which satisfy the same upper bound.
Summation yields
\begin{align}
 \|h_m\|_{C^j}
 &\le C_j\left(1+\sum_{i=1}^m
                  \beta^{i-1}r^{i(1-j)}\right)\notag\\
 &\le C'_j\exp\left(
   [\log\beta+(j-1)_+|\log r|+1]m\right)
 \le C'_j e^{C_*(j+1)m},
 \label{eq:all-order-envelope-growth}
\end{align}
where $(j-1)_+=\max\{j-1,0\}$ and one finite $C_*$ works for
every $j\ge0$.

The lower Hessian margin has the same exponential control.
At a point outside $P_0$, the seed margin remains unchanged.  At a
point in $P_0$, consider the last generation $i\le m$ that changes
the function there.  If no such generation occurs, use the seed
margin.  Otherwise Lemma~\ref{lem:corrugation} bounds its retained
margin below by
\[
 c\,\frac{s_{\min}\alpha^{i-1}}{r^{i-1}}
 \ge c s_{\min}\alpha^m.
\]
Later perturbations vanish near that point, or all their jets there
vanish at their support boundaries, so this margin persists.  By
increasing $C_*$ if necessary,
\begin{equation}\label{eq:iterated-hessian-margin}
 \max_eJ_{h_m}(e)\ge c e^{-C_*m}
 \quad\text{on }\{t\ge1/4\}.
\end{equation}
This argument requires no nonzero-gradient assumption outside
$P_0$; critical points there are covered by the seed lemma.

Fix $\theta>0$ so small that $C_*\theta\le\delta$, and for each
$k\ge2$ take
\[
 m=\lfloor\theta\log k\rfloor,
 \qquad \gamma=\theta\log(1+\zeta)>0,
 \qquad h=h_m+c_k.
\]
Equations \eqref{eq:all-order-envelope-growth} and
\eqref{eq:iterated-hessian-margin} give the first two estimates in
\eqref{eq:envelope-bounds}.  Equation~\eqref{eq:iterated-mass}
gives the third, since
$(1+\zeta)^m\ge(1+\zeta)^{-1}k^\gamma$.
Adding bounded constants affects only the zeroth-order bound, and
$h=-a_*t+c_k$ on $t\le1/2$ because both the seed perturbation and
all corrugations lie above that set.

For clarity about subsequent uses, the constants determining
$C_*$, $\theta$, and $\gamma$ use only the fixed pattern, the
seed, and common $C^1$ metric bounds.  Higher derivatives of an
individual input metric may enter constants in a later smooth
coordinate change, but cannot alter these exponents.  Indeed, for
any fixed such coordinate change, an order-$j$ derivative of the
pullback is a finite sum of derivatives of $h$ of orders at most
$j$ multiplied by fixed smooth coefficients.  Its growth is still
bounded by $C_jk^{\delta(j+1)}$.  Each finite order may have its
own constant and frequency threshold.  This establishes the common
exponent for arbitrarily high finite accuracy and completes the
proof.
\end{proof}
